\documentclass[10pt,letterpaper]{article}

% ── Packages ──────────────────────────────────────────────────────────────────
\usepackage[top=0.5in,bottom=0.45in,left=0.55in,right=0.55in]{geometry}
\usepackage{booktabs}
\usepackage{array}
\usepackage{xcolor}
\usepackage{colortbl}
\usepackage{graphicx}
\usepackage{tabularx}
\usepackage{multicol}
\usepackage{titlesec}
\usepackage{parskip}
\usepackage{microtype}
\usepackage{helvet}
\renewcommand{\familydefault}{\sfdefault}

% ── Colors ────────────────────────────────────────────────────────────────────
\definecolor{critred}{HTML}{C0392B}
\definecolor{headbg}{HTML}{2C3E50}
\definecolor{rowalt}{HTML}{F5F5F5}
\definecolor{accent}{HTML}{E74C3C}
\definecolor{muted}{HTML}{7F8C8D}

% ── Section styling ───────────────────────────────────────────────────────────
\titleformat{\section}{\normalsize\bfseries\color{headbg}}{}{0em}{}[\color{headbg}\titlerule]
\titlespacing{\section}{0pt}{4pt}{2pt}

\setlength{\parindent}{0pt}
\setlength{\parskip}{2pt}
\pagestyle{empty}

% ══════════════════════════════════════════════════════════════════════════════
\begin{document}

% ── Header ────────────────────────────────────────────────────────────────────
\colorbox{headbg}{%
  \parbox{\dimexpr\textwidth-2\fboxsep}{%
    \vspace{4pt}
    \centering
    {\large\bfseries\color{white} POPULATION HEALTH BRIEF}\\[2pt]
    {\small\color{white!80!headbg}
      Top 10 Highest-Risk U.S.\ Counties \textbar{}
      CDC PLACES 2023 + SVI 2022 \textbar{}
      K-Means Risk Stratification (K\,=\,4)}
    \vspace{4pt}
  }%
}

\vspace{4pt}

% ── Two-column body ───────────────────────────────────────────────────────────
\begin{multicols}{2}

% ── Background ────────────────────────────────────────────────────────────────
\section{Background}

County-level CDC PLACES health estimates (2019–2023) were merged with the
CDC/ATSDR Social Vulnerability Index (SVI) on FIPS code and standardized via
\textit{z}-score scaling. K-Means clustering (K\,=\,4) segmented all \textasciitilde3{,}000
U.S.\ counties into four risk tiers: \textbf{Low}, \textbf{Moderate}, \textbf{High},
and \textbf{Critical}. Counties in the \textbf{Critical} tier ($n = 238$, 10\% of
all counties) were ranked by a composite score averaging normalized
diabetes prevalence, hypertension prevalence, and overall SVI. The 10 highest
composite-scoring counties are presented below.

% ── Key Statistics ────────────────────────────────────────────────────────────
\section{Critical vs.\ National Average (2023)}

\begin{tabularx}{\columnwidth}{X r r}
\toprule
\textbf{Indicator} & \textbf{\color{critred}Critical} & \textbf{National} \\
\midrule
Diabetes prevalence       & 15.8\%  & 11.1\% \\
Hypertension prevalence   & 42.9\%  & 33.9\% \\
Physical inactivity       & 35.8\%  & 27.2\% \\
Obesity prevalence        & 44.1\%  & 38.1\% \\
Current smoking           & 21.1\%  & 16.5\% \\
Food insecurity           & 30.5\%  & 17.3\% \\
Lack of transportation    & 15.6\%  &  9.5\% \\
Housing insecurity        & 23.1\%  & 14.0\% \\
SVI (overall, 0--1)       & 0.90    &  0.48  \\
\bottomrule
\end{tabularx}

\vspace{2pt}

% ── Key Findings ──────────────────────────────────────────────────────────────
\section{Key Findings}

\begin{itemize}\setlength\itemsep{1pt}
  \item The top 10 highest-risk counties are concentrated in the \textbf{Deep
        South} (Mississippi 5, Alabama 3, Louisiana 2), reflecting
        persistent regional disparities in chronic disease burden.
  \item All 10 counties have diabetes rates \textbf{$\geq$19.5\%}, nearly
        \textbf{double} the national average of 11.1\%.
  \item Hypertension exceeds \textbf{49\%} in every top-10 county; Holmes County,
        MS reaches \textbf{53.1\%}.
  \item SVI scores cluster near the ceiling ($\geq$0.88), signaling extreme
        concurrent vulnerability across socioeconomic, housing, and
        transportation domains.
  \item Physical inactivity ($r = 0.87$), lack of transportation ($r = 0.91$),
        and food insecurity ($r = 0.94$) are the strongest SDOH correlates of
        county-level diabetes in 2023 (Pearson $r$, $n = 2{,}299$ counties).
\end{itemize}

\end{multicols}

% ── Top 10 Table ──────────────────────────────────────────────────────────────
\section{Top 10 Highest-Risk Counties}

\renewcommand{\arraystretch}{1.2}
\small
\begin{tabularx}{\textwidth}{@{} r l l r r r r r r r @{}}
\toprule
\textbf{Rank} &
\textbf{County} &
\textbf{State} &
\textbf{Pop.} &
\textbf{Diabetes \%} &
\textbf{Hypert. \%} &
\textbf{SVI} &
\textbf{Phys.\ Inact.\ \%} &
\textbf{Obesity \%} &
\textbf{Smoking \%} \\
\midrule
\rowcolor{critred!18}
1 & Holmes     & Mississippi &  15{,}777 & \textbf{21.0} & 53.1 & 0.992 & 44.4 & 51.5 & 23.5 \\
2 & Humphreys  & Mississippi &   7{,}216 & 20.6 & 52.4 & 0.979 & 46.2 & 53.0 & 24.1 \\
\rowcolor{rowalt}
3 & Greene     & Alabama     &   7{,}341 & \textbf{21.1} & 52.4 & 0.916 & 42.0 & 53.4 & 23.8 \\
4 & Bullock    & Alabama     &   9{,}897 & \textbf{21.3} & 52.4 & 0.899 & 42.7 & 51.0 & 24.9 \\
\rowcolor{rowalt}
5 & E.\ Carroll & Louisiana &   6{,}829 & \textbf{21.3} & 52.5 & 0.878 & 44.1 & 51.7 & 29.1 \\
6 & Perry      & Alabama     &   7{,}738 & 20.5 & 51.9 & 0.921 & 40.3 & 54.0 & 23.3 \\
\rowcolor{rowalt}
7 & Sharkey    & Mississippi &   3{,}336 & 19.6 & 51.2 & 0.984 & 44.3 & 50.5 & 23.5 \\
8 & Tunica     & Mississippi &   9{,}234 & 19.9 & 51.2 & 0.960 & 44.8 & 50.8 & 23.6 \\
\rowcolor{rowalt}
9 & Coahoma    & Mississippi &  20{,}077 & 19.5 & 50.7 & 0.990 & 42.1 & 47.1 & 22.0 \\
10 & Madison   & Louisiana   &   9{,}246 & 19.6 & 49.6 & 1.000 & 42.5 & 50.1 & 26.2 \\
\bottomrule
\end{tabularx}

\vspace{3pt}

% ── Recommendations ───────────────────────────────────────────────────────────
\begin{multicols}{2}
\section{Recommended Interventions}

\begin{itemize}\setlength\itemsep{1pt}
  \item \textbf{Food access:} Expand SNAP outreach and mobile food pantry
        programs; prioritize counties with food insecurity $>$40\%.
  \item \textbf{Transportation:} Fund community health worker home-visit models
        and non-emergency medical transport in counties where transportation
        access gaps exceed 15\%.
  \item \textbf{Preventive care:} Deploy federally qualified health center
        (FQHC) satellites targeting hypertension screening and diabetes
        management in Critical tier counties.
  \item \textbf{Physical activity:} Partner with schools and faith communities
        for low-cost activity programs given $>$40\% physical inactivity rates.
\end{itemize}

\section{Data \& Methods Notes}

{\small\color{muted}
Data: CDC PLACES 2023 release (data year 2023); CDC/ATSDR SVI 2022.
Composite risk score = mean of min--max normalized diabetes prevalence,
hypertension prevalence, and overall SVI (RPL\_THEMES).
K-Means validated via elbow and silhouette analysis; K\,=\,4 selected.
Age-adjusted prevalence used where available; crude otherwise.
657 counties excluded from clustering due to missing values in one or more
clustering features (complete-case analysis; 2{,}956 in panel $\to$ 2{,}299 clustered).
}

\end{multicols}

\vspace{1pt}
{\footnotesize\color{muted}
\textit{Prepared:} March 2026 \quad
\textit{Sources:} CDC PLACES (\texttt{data.cdc.gov}), CDC/ATSDR SVI (\texttt{www.atsdr.cdc.gov/placeandhealth/svi}) \quad
\textit{Analysis:} Python 3 · pandas · scikit-learn · K-Means (K\,=\,4)
}

\end{document}
